
% SGA 3，报告 VII_B，英文重建本第一轮的中文译本。
% 范围：依据文献局部页 56--57，第 2.5.1 节全文，止于 2.5.2 前。
% 合并阅读本页 569--570；重版印刷页 558--559；原版边码 523。
% 编者注 97--99。依据：Polo--Gille Exp7B-13oct24.pdf。
% SHA-256 D303690E0195654BBE7C3DAA5D89EA23E24CD709A3575217234AC19AE8836694。
% 参照了 Jacob Reinhold 的 CC BY 4.0 Markdown，作为对照来源。
% 原英文 Loop 1 阶段说明：三幅图均为第一轮直接从 PDF 截取的原图。（当前中文工程已采用原生图。）
\subsubsection*{2.5.1}
\phantomsection
\label{sga3:VIIB:2.5.1}
仍设 \(k\) 是任意伪紧环。由于 \SGARefLink{sga3:VIIB:1.6.F}{1.6.F} 中的函子
\[
  X\longmapsto X_e
\]
与有限积交换，它将每个形式群 \(G\) 映为平展形式群 \(G_e\)。该结论中的态射 \(p:X\to X_e\) 关于 \(X\) 是函子性的，所以此时
\[
  p:G\longrightarrow G_e
\]
是形式群态射。
考察其核 \(\ker(p)\)。\footnote{%
\label{sga3:VIIB:nde:97}以下原论证的细节由编者补充。}它是下图的纤维积
\begin{center}
  \phantomsection
  \label{sga3:VIIB:2.5.1:kernel-fiber-square}
  % SGA3-VII-DIAG-113
  % SGA3-VII-DIAG-113
% Authority: Exp7B-13oct24.pdf, local p. 56.
% Crop: figures/Exp7B_localp056_kernel_etale_fiber_square.png
% Locator: tex/components/48_expose_VIIB_section251_infinitesimal_identity_component_en.tex, Loop 1 line 37.
\begin{tikzcd}[row sep=2.5em,column sep=3.8em]
{} & G \arrow[d,"p"] \\
\operatorname{Spf}(k) \arrow[r,"\varepsilon"]
  & G_e
\end{tikzcd}

\end{center}
其中 \(\varepsilon\) 是 \(G_e\) 的单位截面。由于 \(p\) 在底集合上诱导双射，故\textup{（参见\ \SGARefLink{sga3:VIIB:1.2}{1.2}，编者注 41）}\(\ker(p)\) 的底集合是 \(\operatorname{Spf}(k)\) 在 \(\varepsilon\) 下的像，且对 \(\operatorname{Spf}(k)\) 的每个点 \(s\)，\(\ker(p)\) 在 \(\varepsilon(s)\) 处的局部代数为
\[
  \mathcal O_{G,\varepsilon(s)}
  \widehat\otimes_{\mathcal O_{G_e,\varepsilon(s)}}
  \kappa(s).
\]
在每个点 \(\varepsilon(s)\) 处，\(\mathcal O_{G,\varepsilon(s)}\) 的剩余域为 \(\kappa(s)\)，因而
\[
  \mathcal O_{G_e,\varepsilon(s)}=\kappa(s),
  \qquad
  \mathcal O_{\ker(p),\varepsilon(s)}
  =
  \mathcal O_{G,\varepsilon(s)}.
\]
因此，我们称 \(\ker(p)\) 为 \(G\) 的\emph{单位元的无穷小邻域}，记为
\[
  \ker(p)=G_0.
\]
于是得到正合列
\begin{center}
  \phantomsection
  \label{sga3:VIIB:2.5.1:infinitesimal-exact-sequence}
  % SGA3-VII-DIAG-114
  % SGA3-VII-DIAG-114
% Authority: Exp7B-13oct24.pdf, local p. 57.
% Crop: figures/Exp7B_localp057_infinitesimal_exact_sequence.png
% Locator: tex/components/48_expose_VIIB_section251_infinitesimal_identity_component_en.tex, Loop 1 line 71.
\(
1\longrightarrow G_0\xrightarrow{\ \text{包含}\ }G
\xrightarrow{\ p\ }G_e.
\)

\end{center}
以下，若
\[
  G=G_0.
\]
则称 \(G\) 是\emph{无穷小的}。\footnote{\label{sga3:VIIB:nde:98}下句由编者添加。}
这等价于说，对每个 \(s\in\operatorname{Spf}(k)\)，或者等价地，对每个连续态射 \(k\to\kappa\)（其中 \(\kappa\) 是赋予离散拓扑的域），单位元是 \(G(\kappa(s))\)（分别是 \(G(\kappa)\)）中的唯一元素。
再设 \(G\) 在 \(k\) 上拓扑平坦。\footnote{\label{sga3:VIIB:nde:99}以下细节由编者补充。}由 \SGARefLink{sga3:VIIB:1.6.1}{1.6.1}，态射
\[
  p:G\longrightarrow G_e
\]
是拓扑平坦的。由于它是双射，\SGARefLink{sga3:VIIB:1.3.1}{命题 1.3.1} 表明它是有效满态射。因此 \(G_e\) 可与商 \(G/G_0\) 等同。我们证明了：
\medskip
\noindent\textbf{推论。}
\phantomsection
\label{sga3:VIIB:2.5.1:corollary}
{\itshape 设 \(G\) 是在 \(k\) 上拓扑平坦的形式群。则 \(G_e\) 可与商 \(G/G_0\) 等同；换言之，有形式群的正合列}
\begin{center}
  \phantomsection
  \label{sga3:VIIB:2.5.1:short-exact-sequence}
  % SGA3-VII-DIAG-115
  % SGA3-VII-DIAG-115
% Authority: Exp7B-13oct24.pdf, local p. 57.
% Crop: figures/Exp7B_localp057_infinitesimal_etale_short_exact_sequence.png
% Locator: tex/components/48_expose_VIIB_section251_infinitesimal_identity_component_en.tex, Loop 1 line 107.
\(
1\longrightarrow G_0\xrightarrow{\ \text{包含}\ }G
\xrightarrow{\ p\ }G_e\longrightarrow1.
\)

\end{center}
